\section{Introduction}\label{s:intro}
Many OEIS entries end with a line beginning ``Conjecture:''. The statement was usually found by computing
terms, and nobody has written a proof or a refutation since. This report collects such statements that we
were able to settle. It started as a three-page note on nine entries about binary morphic words
(Part~\ref{part:one}, reproduced with minor changes), and grew as the same methods, and others, were applied to
a systematic crawl of the OEIS.

\subsection{What is counted}
An entry is counted once, when at least one conjecture written in it (or a formula marked ``conjectured'')
is settled here and we found no earlier proof or refutation. Statements that appear without the word
``conjecture'' (``it appears that'', or a limit stated without proof) are counted only when they are worded
as conjectures; the $25$ further complementary equations of Section~\ref{s:compeq} are proved but not
counted. Entries already settled by others (Dekking, Zabolotskii, Shallit, Bosma et al., Ashburn, Song and
others) are listed in Part~\ref{part:five} and are not counted.

\begin{center}
\begin{tabular}{lr}\toprule
Result & Entries\\\midrule
true as stated (possibly after an obvious reading of the index) & \Ntrue\\
false as stated & \Nfalse\\
mixed: one statement true, another false & \Nmixed\\
true for $n\ge2$ only & \Ntwo\\\midrule
total & \Ntotal\\\bottomrule
\end{tabular}
\end{center}
The complete list, with the author of each conjecture, the result, the method and the section with the
proof, is Appendix~\ref{a:catalogue}.

\subsection{The main results}
The results differ a lot in depth. We list first those that we believe need a real argument.
\begin{itemize}
\item \textbf{Real roots of the Doche--Mend\`es France polynomials} (\oeis{A347171},
  Theorem~\ref{t:A347171}). $P_n(y)=\sum_{j<2^n}\mathrm{GRS}(j)\,y^{\mathrm{wt}(j)}$ has exactly
  \oeis{A285869}$(n)$ real roots, all simple, with an explicit count in each of six regions.
\item \textbf{Hadjicostas' Josephus conjecture for every $k\ge2$} (\oeis{A005427}, \oeis{A120160} and the
  entries for $k=2,\dots,9$; Theorem~\ref{t:josephus}), with the explicit survivor $J_k(T(n))=r_n+1$. It turns
  out to be a short corollary of the Odlyzko--Wilf formula.
\item \textbf{Kotesovec's exponent} for \oeis{A371176}: $\log a(n)/\log n\to\log2/\log\varphi$, through a
  Fibonacci count of Schreier sets (Theorem~\ref{t:A371176}).
\item \textbf{Gutkovskiy's nested recursions} $a_k(n)=n-a_k(\lfloor a_k(n-1)/k\rfloor)$ for every $k\ge1$
  (\oeis{A366860}, \oeis{A366870}, \oeis{A366871}, \oeis{A368281}): $a_k$ is the counting function of the
  fixed point of $0\mapsto1^k$, $1\mapsto1^k0$ (Theorem~\ref{t:A366860}).
\item \textbf{Kimberling's discrepancy conjectures} for binary morphic words: exact extreme values of
  $nr-a(n)$ in $\mathbb{Q}(\sqrt d)$ (Parts~\ref{part:one} and~\ref{part:two}). The smallest counterexample to
  \oeis{A284364} is $n=2\,977\,771$.
\item \textbf{Stephan's conjecture for \oeis{A005590}}: $a(3n)=0$ if and only if $n$ is fibbinary, by an
  invariant quadratic cone (Theorem~\ref{t:A005590}).
\item \textbf{Cloitre's constant for binary partitions is wrong} (\oeis{A033485}): the limit of
  $a(2n)/a(n)\cdot\log n/n$ exists and is $2\log2=1.386\ldots$, not $1.64\ldots$ (Theorem~\ref{t:A033485}); this also follows from
  de Bruijn's 1948 asymptotics.
\item \textbf{Song's lower bound for Golomb's sequence}, $a(n)\ge n^{\varphi-1}$ (Theorem~\ref{t:A001462}).
\item \textbf{The anti-tetranacci system} (\oeis{A299405}): the published bound for $c$ fails at
  $c(26)=113$, infinitely often (Section~\ref{s:A299405}). This entry is not counted, because published work
  treats the system; we could not read it, and the disagreement should be checked there.
\end{itemize}
Many other results are short. Some are classical arguments that had not been written down in the entry
(Shevelev's conjecture in \oeis{A014664} is the lifting-the-exponent lemma). We count them all, since each
entry still presented the statement as open, but we mark the trivial ones as such.

\subsection{Methods}
\begin{description}
\item[Hand proofs.] Most entries in Parts~\ref{part:three} and~\ref{part:four}.
\item[\Walnut.] An automatic theorem prover for first-order statements about automatic sequences, by
  H.~Mousavi, developed by J.~Shallit and others~\cite{Mo,Sh}. A \TRUE{} answer is a proof for all $n$. We
  use \Walnut{} (commit \texttt{a4fe84e}) in base $2$, $3$, $4$, Fibonacci numeration, and Dumont--Thomas
  numerations of non-uniform morphisms, the latter validated inside \Walnut{} by axioms. Every script
  contains at least one statement that must be \FALSE, as a control.
\item[Exact discrepancy.] The supremum and infimum of $nr-a(n)$ for a fixed point of a binary Pisot morphism
  are the solution of a contracting max/min system, solved exactly in $\mathbb{Q}(\sqrt d)$ with a checked
  certificate (Section~\ref{s:method}).
\item[Perron frequencies.] For primitive substitutions, letter and factor frequencies are the normalised
  Perron eigenvectors (Queffélec~\cite{Qu}). This settles limit conjectures $a(n)/n\to m$.
\item[Exact algebra.] For algebraic generating functions, a conjectured holonomic recurrence is reduced
  exactly modulo the minimal polynomial.
\item[Published theorems.] Denisov--Wachtel for walks in cones (\oeis{A361190}) and
  Alekseyev--Amdeberhan--Shallit--Vukusic for $\nu_3$ of Delannoy numbers (\oeis{A358360}).
\end{description}

\subsection{Reliability}
None of these proofs has been refereed. The work was done with the help of an AI system; every proof was
checked by programs written from the OEIS definitions, independently of the code that found the result.
On November 19, 2026, the whole verification suite was re-run from a fresh clone of the repository: all
seventeen \Walnut{} scripts give exactly the recorded answers, and the Python checks pass (Appendix~\ref{a:verif}).
Each falsity claim was also re-checked against the data of the entry fetched from the OEIS on that day.
The search for prior work was a web and arXiv search and a reading of each entry; for some entries the
original paper could not be read, and we say so. Before any of these results is submitted to the OEIS,
the corresponding prior-work check should be completed.

\paragraph{Conventions.} $\varphi=(1+\sqrt5)/2$; $t(n)$ is the Thue--Morse sequence; $s(n)$ or
$\mathrm{wt}(n)$ is the binary weight; $\nu_p$ is the $p$-adic valuation. For a word $x$ and a letter $L$,
$a(n)$ is the position, counted from $1$, of the $n$-th occurrence of $L$. File names refer to the
repository \url{https://github.com/joaquimmj/oeis-walnut}, branch \texttt{conjectures}; each entry has a
folder \texttt{conjectures/Axxxxxx} with the proof, the checks and a draft OEIS comment.
